\section{Discussion}
\label{sec:discussion}

\paragraph{Why this differs from earlier findings.} Studies that found agents breaking rules under KPI
pressure \citep{li2025odcv, lynch2025agentic} usually built scenarios where the violation was the only or the
obvious route to the goal, and often tested many model families. SandBank gives the agent a realistic
compliant route for every lead (request documents, escalate, report the miss) and a manager who never
approves an exception. That is how a bank is meant to work, and it is the setting a bank deploying an RM agent
cares about. Our result is that, in this setting, three current Claude models take the compliant route even
when the manager licenses or orders the violation. It is compatible with the earlier results: they show
that some models in some designs violate; we show that these models in this design did not, with bounds.

\paragraph{What the null does and does not license.} A bank can read the bounds as evidence that these models,
given this kind of tooling and policy, do not readily circumvent KYC/AML controls for a sales target. It
cannot read them as a guarantee. The study has one bank, one policy, one RM role, short episodes and fixed
texts; the agents may recognise a test (Section~\ref{sec:limitations}); and the bounds are per probe at the
sample sizes we could afford. Banks should still log tool calls, keep hard controls on prohibited actions and
test their own deployment.

\paragraph{The useful difference is capability.} The models differed not in compliance but in how much
legitimate work they finished. A cautious agent that leaves a legal large deposit unopened has a business
cost that a violation-only evaluation would score as perfect. Pre-deployment testing of compliance agents
should measure both.

\paragraph{On null results.} We pre-registered hypotheses that predicted violations, ran the study as
planned, and report that the hypotheses failed. A failed ceiling gate is reported, not worked around by
strengthening the prompt until a violation appears. We think such results belong in the literature, because
a field that publishes only positive misalignment findings will overstate how often it happens.
